% Options for packages loaded elsewhere
\PassOptionsToPackage{unicode}{hyperref}
\PassOptionsToPackage{hyphens}{url}
\PassOptionsToPackage{dvipsnames,svgnames,x11names}{xcolor}
%
\documentclass[
  11pt,
  a4paper,
]{ltjsarticle}
\usepackage{amsmath,amssymb}
\usepackage{iftex}
\ifPDFTeX
  \usepackage[T1]{fontenc}
  \usepackage[utf8]{inputenc}
  \usepackage{textcomp} % provide euro and other symbols
\else % if luatex or xetex
  \usepackage{unicode-math} % this also loads fontspec
  \defaultfontfeatures{Scale=MatchLowercase}
  \defaultfontfeatures[\rmfamily]{Ligatures=TeX,Scale=1}
\fi
\usepackage{lmodern}
\ifPDFTeX\else
  % xetex/luatex font selection
\fi
% Use upquote if available, for straight quotes in verbatim environments
\IfFileExists{upquote.sty}{\usepackage{upquote}}{}
\IfFileExists{microtype.sty}{% use microtype if available
  \usepackage[]{microtype}
  \UseMicrotypeSet[protrusion]{basicmath} % disable protrusion for tt fonts
}{}
\makeatletter
\@ifundefined{KOMAClassName}{% if non-KOMA class
  \IfFileExists{parskip.sty}{%
    \usepackage{parskip}
  }{% else
    \setlength{\parindent}{0pt}
    \setlength{\parskip}{6pt plus 2pt minus 1pt}}
}{% if KOMA class
  \KOMAoptions{parskip=half}}
\makeatother
\usepackage{xcolor}
\usepackage[margin=24mm]{geometry}
\setlength{\emergencystretch}{3em} % prevent overfull lines
\providecommand{\tightlist}{%
  \setlength{\itemsep}{0pt}\setlength{\parskip}{0pt}}
\setcounter{secnumdepth}{5}
\ifLuaTeX
  \usepackage{selnolig}  % disable illegal ligatures
\fi
\IfFileExists{bookmark.sty}{\usepackage{bookmark}}{\usepackage{hyperref}}
\IfFileExists{xurl.sty}{\usepackage{xurl}}{} % add URL line breaks if available
\urlstyle{same}
\hypersetup{
  colorlinks=true,
  linkcolor={blue},
  filecolor={Maroon},
  citecolor={Blue},
  urlcolor={blue},
  pdfcreator={LaTeX via pandoc}}

\author{}
\date{}

\begin{document}

\hypertarget{ux30d9ux30afux30c8ux30eb}{%
\section{ベクトル}\label{ux30d9ux30afux30c8ux30eb}}

ベクトルはしばしば「向きと大きさを持つ矢印」と説明されます。しかし線形代数では、より本質的に
\textbf{加法とスカラー倍ができる対象} と考えます。

\[
x+y,\qquad \alpha x
\]

が同じ種類の対象として定義できることが重要です。

\hypertarget{ux5ea7ux6a19ux306fux30d9ux30afux30c8ux30ebux305dux306eux3082ux306eux3067ux306fux306aux3044}{%
\subsection{座標はベクトルそのものではない}\label{ux5ea7ux6a19ux306fux30d9ux30afux30c8ux30ebux305dux306eux3082ux306eux3067ux306fux306aux3044}}

\[
x=\begin{bmatrix}2\\3\end{bmatrix}
\]

という数の並びは、ある基底を選んだときの座標です。基底を変えれば同じ幾何学的ベクトルでも座標は変わります。

この区別は後の

\begin{itemize}
\tightlist
\item
  基底変換
\item
  固有ベクトル
\item
  PCA
\item
  SVD
\end{itemize}

に直結します。

\hypertarget{ux30d9ux30afux30c8ux30ebux3068ux3057ux3066ux6271ux3048ux308bux3082ux306e}{%
\subsection{ベクトルとして扱えるもの}\label{ux30d9ux30afux30c8ux30ebux3068ux3057ux3066ux6271ux3048ux308bux3082ux306e}}

ベクトルは数の列に限りません。

\begin{itemize}
\tightlist
\item
  多項式
\item
  関数
\item
  画像
\item
  音声波形
\item
  行列
\end{itemize}

なども、加法とスカラー倍についてベクトル空間を作れます。

例えば画像を \texttt{m×n} 個の画素値として見れば、一枚の画像は
\texttt{R\^{}\{mn\}} のベクトルと考えられます。この視点が PCA
や画像圧縮につながります。

\hypertarget{ux7406ux89e3ux30c1ux30a7ux30c3ux30af}{%
\subsection{理解チェック}\label{ux7406ux89e3ux30c1ux30a7ux30c3ux30af}}

「ベクトル」と「その座標」の違いを説明できますか。これを説明できることが、後の基底概念の出発点です。

\end{document}
